We agree that $H(\mathbf{l})$ centers the tendon lengths in their actuation range and does not measure effort, which was not recorded. We now describe it as actuator-range centering and removed the statements about effort (Introduction, Section~3.3.2, Conclusion). To isolate its effect without new hardware experiments, we performed offline ablations on identical targets for both controllers (new SI Section~5.3, Fig.~S8, and Table~S13). For the model-based controller, replaying the control law on the PCC model for the 20 evaluation targets gives a final position error of $5.25 \pm 2.23$~mm with and $5.84 \pm 1.85$~mm without the null-space term, while the RMS deviation of the tendon lengths from $\bar{\mathbf{l}}$ decreases from $12.00 \pm 3.18$~mm to $7.48 \pm 1.96$~mm ($H(\mathbf{l})$: $693 \pm 396$ vs.\ $269 \pm 135$). For the data-driven controller, we retrained the inverse-kinematics network with and without the secondary term ($\zeta = 0.01$ and $\zeta = 0$, same data, split, and seed) and evaluated it through the frozen forward network on the same targets: the secondary term reduces $H(\mathbf{l})$ from $487 \pm 257$ to $371 \pm 166$ (RMS tendon deviation from 10.1 to 8.8~mm), at the cost of a small increase of the position error predicted by the forward network ($5.1 \pm 2.1$ vs.\ $4.6 \pm 2.3$~mm). Unlike the null-space projection, the penalty in the training loss trades some primary-task accuracy for range centering, which we now state in Section~3.3.2. These results are simulated through the respective models and are reported as such.
